\chapter{\MakeUppercase{Conclusion and Further enhancements}}
\justifying
\section{Conclusion}
The aim of the project was to assess whether or not the sensor noise from cameras and microphones could create a practical device verification prototype. The answer is yes within the bounded and defined limits of the results of the system designed and the artefacts created by the system. QNA-Auth demonstrated that multi-source sensor-noise recordings can provide canonical features that can be constructed into a learned similarity space and combined together with other processes into a service-oriented workflow (enrollment, verification, replay-resistance, and operational explainable).

The single best feature of the system is not one metric; it is the combination of:
\begin{itemize}
	\item a complete end-to-end software stack,
	\item Explicit claim boundaries,
	\item Reusable canonical feature pipeline,
	\item Confidence-band based runtime decisions,
	\item Challenge-response protocol that is security aware.
	\item Documentation and artifacts that permit academic peer review.
\end{itemize}

The academic value of the project is established because it goes beyond the confines of an artificially narrow experiment and demonstrates the type of engineering integration expected in a solid capstone thesis.

\section{Further Enhancements}
There are many extensions to the project that would add significant value, including, but not limited to: 
\begin{itemize}
	\item Collecting larger multi-session datasets, using more diverse devices and environments.
	\item Performing matched ablation testing on camera only metrics, microphone only metrics, and the fused configuration.
	\item Calibrating thresholds based upon the environment in which QNA-Auth will be used, instead of using a single global setting.
	\item Implementing more robust anti-spoofing and liveness checks.
	\item Developing a strong client-assisted possession model in the challenge verification process.
	\item Providing more granular longitudinal drift analysis.
	\item Established formal reporting-ready statistical confidence intervals based upon larger evaluation runs.
\end{itemize}

In effect, QNA-Auth should be presented as a security-aware prototype for high-confidence device matching, not as a fully developed identity platform. With that in mind, it is both technically sound and thesis-worthy.